% ===== PRÉAMBULE CANONIQUE — à coller VERBATIM en tête de CHAQUE .tex =====
\documentclass[11pt,a4paper]{article}
\usepackage[utf8]{inputenc}\usepackage[T1]{fontenc}\usepackage{lmodern}
\usepackage[french]{babel}
\usepackage{amsmath,amssymb,mathtools}\usepackage{icomma}
\usepackage[version=4]{mhchem}          % \ce{...} : équations & formules
\usepackage{siunitx}\sisetup{locale=FR} % \qty{}{}, \si{}, \ang{}
\usepackage{chemfig}                    % \chemfig{...} : structures développées
\usepackage{xcolor}\usepackage{colortbl}
\usepackage{tikz}\usepackage{pgfplots}\pgfplotsset{compat=1.18}
\usepackage[most]{tcolorbox}\usepackage{enumitem}\usepackage{array,booktabs}
\usepackage[a4paper,margin=2.2cm]{geometry}
\usetikzlibrary{arrows.meta,calc,angles,quotes,positioning,patterns,backgrounds,shapes.geometric}
\definecolor{primary}{RGB}{222,49,99}\definecolor{primarydark}{RGB}{150,22,55}
\definecolor{defcol}{RGB}{0,114,178}\definecolor{methcol}{RGB}{52,92,148}
\definecolor{excol}{RGB}{0,135,90}\definecolor{attncol}{RGB}{200,40,55}
\tcbset{boxbase/.style={enhanced,breakable,boxrule=0.9pt,arc=2.5pt,
  left=3mm,right=3mm,top=2.3mm,bottom=2.3mm,fonttitle=\bfseries,coltitle=white}}
% --- boîtes TUTORIEL (noms EXACTS, ne pas renommer) ---
\newtcolorbox{cle}[1][]{boxbase,colback=defcol!6,colframe=defcol,
  title={\textsc{À retenir}\ifx\relax#1\relax\relax\else\textmd{ : \itshape #1}\fi}}
\newtcolorbox{methode}[1][]{boxbase,colback=methcol!7,colframe=methcol,sharp corners,
  title={\textsc{Méthode}\ifx\relax#1\relax\relax\else\textmd{ --- \itshape #1}\fi}}
\newtcolorbox{exemplebox}[1][]{boxbase,colback=excol!5,colframe=excol,
  title={\textsc{Exemple}\ifx\relax#1\relax\relax\else\textmd{ : \itshape #1}\fi}}
\newtcolorbox{piege}[1][]{boxbase,colback=attncol!6,colframe=attncol,
  title={\textsc{Piège}\ifx\relax#1\relax\relax\else\textmd{ : \itshape #1}\fi}}
\newtcolorbox{remarque}[1][]{boxbase,colback=black!4,colframe=black!45,
  title={\textsc{Remarque}\ifx\relax#1\relax\relax\else\textmd{ : \itshape #1}\fi}}
% --- boîtes EXERCICE / CORRIGÉ (noms EXACTS, auto-comptées) ---
\newtcolorbox[auto counter]{exo}[1][]{boxbase,colback=primary!4,colframe=primary,
  title={\textsc{Exercice}~\thetcbcounter\ifx\relax#1\relax\relax\else\textmd{ --- \itshape #1}\fi}}
\newtcolorbox[auto counter]{corrige}[1][]{boxbase,colback=excol!4,colframe=excol,
  title={\textsc{Corrigé}~\thetcbcounter\ifx\relax#1\relax\relax\else\textmd{ --- \itshape #1}\fi}}
\newcommand{\R}{\mathbb{R}}\newcommand{\N}{\mathbb{N}}\newcommand{\Z}{\mathbb{Z}}\newcommand{\ds}{\displaystyle}
% (un \usepackage{fancyhdr}+en-tête est possible ; il est ignoré côté web)
% =========================================================================
\begin{document}

\begin{corrige}[lire $Z$ et $A$ dans la notation]
Le nombre du bas est $Z=n(p^+)$ ; celui du haut est $A$, le nombre de nucléons.
\begin{enumerate}[label=\alph*)]
  \item $\ce{^{19}_{9}F}$ : $Z=9$, $A=19$, donc $n(p^+)=9$ protons et $19$ nucléons.
  \item $\ce{^{31}_{15}P}$ : $Z=15$, $A=31$, donc $n(p^+)=15$ protons et $31$ nucléons.
  \item $\ce{^{7}_{3}Li}$ : $Z=3$, $A=7$, donc $n(p^+)=3$ protons et $7$ nucléons.
\end{enumerate}
\end{corrige}

\begin{corrige}[compter les neutrons]
On applique $n(n^0) = A - Z$.
\begin{enumerate}[label=\alph*)]
  \item $\ce{^{14}_{7}N}$ : $n(n^0) = 14 - 7 = 7$ neutrons.
  \item $\ce{^{23}_{11}Na}$ : $n(n^0) = 23 - 11 = 12$ neutrons.
  \item $\ce{^{40}_{18}Ar}$ : $n(n^0) = 40 - 18 = 22$ neutrons.
\end{enumerate}
\end{corrige}

\begin{corrige}[électrons d'un atome neutre]
Atome neutre $\Rightarrow n(e^-) = n(p^+) = Z$.
\begin{enumerate}[label=\alph*)]
  \item $\ce{^{12}_{6}C}$ : $n(e^-) = 6$ électrons.
  \item $\ce{^{27}_{13}Al}$ : $n(e^-) = 13$ électrons.
  \item $\ce{^{32}_{16}S}$ : $n(e^-) = 16$ électrons.
\end{enumerate}
\end{corrige}

\begin{corrige}[écrire la notation d'un noyau]
$Z = n(p^+)$ (en bas) et $A = n(p^+) + n(n^0)$ (en haut) ; le nombre de nucléons est $A$.
\begin{enumerate}[label=\alph*)]
  \item $Z=17$, $A = 17+18 = 35$ : $\ce{^{35}_{17}Cl}$, soit $35$ nucléons.
  \item $Z=11$, $A = 11+12 = 23$ : $\ce{^{23}_{11}Na}$, soit $23$ nucléons.
  \item $Z=8$, $A = 8+8 = 16$ : $\ce{^{16}_{8}O}$, soit $16$ nucléons.
\end{enumerate}
\end{corrige}

\begin{corrige}[charge d'un ensemble de particules]
La charge totale est le nombre de particules multiplié par la charge de chacune ($+e$ pour un proton,
$-e$ pour un électron), avec $e = \SI{1,6e-19}{\coulomb}$.
\begin{enumerate}[label=\alph*)]
  \item $Q = +e = +\SI{1,6e-19}{\coulomb}$.
  \item $Q = 4 \times (+e) = +4 \times \SI{1,6e-19}{\coulomb} = +\SI{6,4e-19}{\coulomb}$.
  \item $Q = 5 \times (-e) = -5 \times \SI{1,6e-19}{\coulomb} = -\SI{8,0e-19}{\coulomb}$.
\end{enumerate}
\end{corrige}

\end{document}
